\begin{abstract}
Under calibrated public density envelopes satisfying \(c_{\mathrm{lo}}\le(\kappa+1)2^\kappa\le c_{\mathrm{hi}}\), where \(c_{\mathrm{lo}}\) and \(c_{\mathrm{hi}}\) are the density-envelope constants and \(\kappa\) is the design-decay exponent, we characterize optimal estimation and honest interval length for the population causal mean at a prespecified latent dose when treatment is measured with Gaussian error and local dose support may vanish polynomially. The model has two observed strata, compact latent-dose support, bounded potential outcomes, schedule-level conditional exchangeability, and classical measurement error with known standard deviation \(\sigma\in[0,1/4]\). Conditional mean functions have Hölder exponent \(\beta\in(0,1]\), and conditional dose densities satisfy these public polynomial envelopes with \(\kappa\in[0,2]\). For this calibrated class and any fixed noncoverage probability \(\alpha\in(0,1/2)\), minimax expected absolute error and minimum worst-law expected length of honest connected \((1-\alpha)\)-coverage intervals have matching orders, uniformly over the error scale. For sample size \(n\), the characterization combines a direct weak-design regime with rate \(n^{-\beta/(2\beta+\kappa+1)}\), an intermediate Gaussian-inversion regime, and a compact-support regime. Neighboring rate formulas are comparable throughout fixed multiplicative neighborhoods of both transition thresholds. Every fixed error scale \(\sigma\in(0,1/4]\) yields order \((\log\log n/\log n)^\beta\) for this calibrated class. An observable estimator selects among Fourier, polynomial, and constant inverse weights using public error certificates, with the declared small-sample fallback; its associated interval has finite-sample coverage \(1-\alpha\) for every \(\alpha\in(0,1)\). Under the same envelope calibration, matching lower bounds follow from bounded Bernoulli structural alternatives whose causal means are separated while their complete observed samples remain difficult to distinguish.
\end{abstract}

\section{Introduction}

This paper characterizes how measurement precision and local treatment support jointly determine the accuracy of pointwise causal estimation under calibrated public density envelopes. We study the population mean response to setting a latent continuous treatment to the prespecified dose \(1/2\). Treatment is observed with independent Gaussian measurement error of known standard deviation \(\sigma\), the error scale. Within the bounded structural law class in \cref{obj:def:model-class}, assuming \(c_{\mathrm{lo}}\le(\kappa+1)2^\kappa\le c_{\mathrm{hi}}\), we establish matching orders for minimax expected absolute error and minimum worst-law expected length of honest connected \((1-\alpha)\)-coverage intervals for \(\alpha\in(0,1/2)\). This calibrated-class characterization is uniform over the maintained range \(\sigma\in[0,1/4]\), including error-free observation, for fixed public class constants and exponents in their declared ranges. An observable estimator and interval attain these orders, and structural alternatives evaluated through their induced observed distributions supply the converse under the same envelope calibration.

The model combines causal adjustment with explicit restrictions on dose availability and response regularity. Observations consist of a stratum indicator taking two values, a contaminated treatment measurement, and a realized outcome. The latent dose lies in \([0,1]\). Public constants \(K=(p_{\min},c_{\mathrm{lo}},c_{\mathrm{hi}})\) use \(p_{\min}\) as the minimum stratum probability and \(c_{\mathrm{lo}},c_{\mathrm{hi}}\) as the density-envelope constants. Thus stratum probabilities lie in \([p_{\min},1-p_{\min}]\), while the density of the centered latent-dose displacement \(t=A-1/2\) lies between \(c_{\mathrm{lo}}|t|^\kappa\) and \(c_{\mathrm{hi}}|t|^\kappa\), where \(0<p_{\min}\le1/2\) and \(0<c_{\mathrm{lo}}\le c_{\mathrm{hi}}\). Conditional potential-outcome means satisfy radius-one Hölder increment bounds with exponent \(\beta\in(0,1]\), the mean-smoothness exponent, and take values in \([0,1]\); potential outcomes are bounded in the same interval. We maintain \(\kappa\in[0,2]\). Conditional exchangeability applies to the entire measurable potential-outcome schedule, consistency links its random-dose evaluation to the realized outcome, and the Gaussian error is independent of the structural coordinates. Under these conditions, \cref{obj:prop:identification} identifies the population causal mean even when the dose-density envelopes vanish at the evaluation point.

The main result describes three ranges of measurement precision under the envelope calibration \(c_{\mathrm{lo}}\le(\kappa+1)2^\kappa\le c_{\mathrm{hi}}\). Let \(n\) denote sample size, let \(d=2\beta+\kappa+1\) be the effective smoothness-and-design exponent, and let \(L_n=\log(en)\) be the logarithmic sample-size scale. The localization scale \(b_{n,\sigma}\) is
\[
b_{n,\sigma}=
\begin{cases}
n^{-1/d},
&0\le\sigma\le n^{-1/d},\\[3pt]
\dfrac{\sigma}{\sqrt{\log(e+n\sigma^d)}},
&n^{-1/d}<\sigma\le L_n^{-1/2},\\[7pt]
\dfrac{\log(e+\sigma^2L_n)}{L_n},
&L_n^{-1/2}<\sigma\le1/4.
\end{cases}
\]
For this calibrated class and all sufficiently large samples, \cref{obj:thm:uniform-frontier} bounds both optimal criteria above and below by constant multiples of \(b_{n,\sigma}^{\beta}\), uniformly over the known error scale and the declared public potential-outcome spaces. The direct branch balances mean smoothness against the scarcity of nearby latent doses. Gaussian Fourier inversion governs the intermediate branch, while known compact support supports polynomial localization in the third branch. For every fixed multiplicative neighborhood of either threshold, the adjacent formulas are comparable once sample size is sufficiently large. Within the calibrated class, at zero error, both criteria have order \(n^{-\beta/(2\beta+\kappa+1)}\); at every fixed positive error scale in the stated range, both eventually have order \((\log\log n/\log n)^\beta\).

The upper bound comes from positive latent localization weights paired with observable inverses of the Gaussian channel. Weighted outcome and treatment moments estimate a localized mean within each stratum, and empirical stratum frequencies aggregate those estimates into the population target. A finite dictionary contains Fourier weights, polynomial inverse-heat weights, and a constant entry. The estimator selects a dictionary entry by minimizing a public expected-error certificate determined by sample size, the known error scale, the public class constant \(K\), and the smoothness and design-decay exponents \(\beta,\kappa\). Thus the construction in \cref{obj:def:total-estimator} uses one selection rule across the three regimes. \Cref{obj:thm:observable-upper} establishes its uniform large-sample risk bound. Dividing the same certificate by \(\alpha\) calibrates a connected interval with coverage at least \(1-\alpha\) for every sample size and every law in the model class, while its expected-length bound attains the frontier for sufficiently large samples, as stated in \cref{obj:thm:finite-honesty}.

Under the envelope calibration \(c_{\mathrm{lo}}\le(\kappa+1)2^\kappa\le c_{\mathrm{hi}}\), the matching lower bound uses Bernoulli threshold schedules within the bounded structural class. The alternatives share their latent treatment design and measurement channel, while opposite perturbations change the causal mean at the evaluation dose. In the inverse regimes, normalized filtered Jacobi packets preserve the perturbation's pointwise height and cancel low-order moments under the polynomial design weight. Their localization and regularity bounds control model membership, and Gaussian convolution attenuates the resulting change in the observed outcome-marked distribution. \Cref{obj:thm:observed-lower} supplies target separation at the frontier order together with a uniform total-variation bound for the complete observed sample under this calibration. This comparison yields lower bounds for both estimation and honest connected-interval length over the same calibrated class used for achievability.

The closest causal measurement-error comparison is the contaminated-treatment dose-response estimator of \citet{HuangZhang2023}. Under the conditions of their arXiv version 2, Theorem 4.4, its fixed-dose expansion has bias orders \(h^2+h_0^2\) and stochastic remainder \(O_p\{N^{-1/2}/\min(e(h),e(h_0))\}\), where \(e(b)=b^{1/2}\exp(-b^{-r}/\gamma)\) in the supersmooth case. Our contribution is a matched minimax characterization of absolute risk and honest interval length under explicit polynomial design envelopes satisfying \(c_{\mathrm{lo}}\le(\kappa+1)2^\kappa\le c_{\mathrm{hi}}\), uniform over the known noise scale. The direct endpoint shares its rate exponent with the degenerate-design Gaussian-response regression benchmark of \citet[Definition 2 and Theorem 1]{Gaiffas2005}, as made precise in \cref{obj:prop:error-free-reduction}. Gaussian attenuation connects the intermediate branch to errors-in-variables regression \citep[Sections 3--4]{FanTruong1993}, while compact-support density deconvolution provides a precedent for iterated-logarithm-over-logarithm localization \citep[Theorems 1--2]{Meister2007}. These comparisons concern their respective experiments, restrictions, and losses. Under the same calibrated envelopes, the present characterization places the three mechanisms within one causal model and couples their estimation orders to honest uncertainty about the same scalar target.



The paper proceeds through related work, setup and identification, and the optimal risk and interval-length characterization. It then presents the observable estimator and honest intervals, followed by the observed-law lower bounds. Discussion and limitations and future work follow; the appendices develop identification and error certificates, weighted packets and the observed-experiment converse, and proofs of the main results.
